% Part: normal-modal-logic
% Chapter: tableaux
% Section: simple-S5

\documentclass[../../../include/open-logic-section]{subfiles}

\begin{document}

\olfileid{nml}{tab}{s5}

\olsection{\usetoken{P}{tableau} مبسطة للنسق \Log{S5}}

النسق \Log{S5} سليم وتام بالنسبة إلى فئة النماذج الكلية، أي النماذج
التي يكون فيها كل عالم متاحًا من كل عالم. ولا تؤثر علاقة الإتاحة في
النماذج الكلية: فالقول «يوجد عالم~$w$ بحيث $\mSat{M}{!A}[w]$» صادق
إذا وفقط إذا وُجد عالم كهذا متاح من~$u$. ومن ثم يمكننا في \Log{S5}
أن نعرّف النموذج بأنه مجرد مجموعة من العوالم وإسناد~$V$. وهذا يوحي
بإمكان تبسيط قواعد ال!!{tableau} أيضًا. ففي الحالة العامة نتخذ
متتاليات الأعداد الصحيحة الموجبة بادئات، كي نتعقب أي تلك البادئات
تسمي عوالم متاحة من غيرها: إذ تسمي $\sigma.n$ عالمًا متاحًا
من~$\sigma$. لكن كل عالم في \Log{S5} متاح من كل عالم، فلا حاجة إلى
هذا التعقب. ويمكننا بدلًا من ذلك استعمال الأعداد الصحيحة الموجبة
بادئات. وترد القواعد المبسطة في \olref{tab:rules-S5}.

\begin{table}
  \begin{center}
    \def\arraystretch{3}\def\fCenter{}
    \begin{tabular}{|c|c|}
    \hline
    \iftag{prvBox}{
      \AxiomC{\sFmla{\True}{\Box !A}[n]}
      \RightLabel{\TRule{\True}{\Box}}
      \UnaryInfC{\sFmla{\True}{!A}[m]}
      \DisplayProof
      &
      \AxiomC{\sFmla{\False}{\Box !A}[n]}
      \RightLabel{\TRule{\False}{\Box}}
      \UnaryInfC{\sFmla{\False}{!A}[m]}
      \DisplayProof\\
      $m$ مستعملة & $m$ جديدة
      \\[1ex]
      \hline}{}
    \iftag{prvDiamond}{
      \AxiomC{\sFmla{\True}{\Diamond !A}[n]}
      \RightLabel{\TRule{\True}{\Diamond}}
      \UnaryInfC{\sFmla{\True}{!A}[m]}
      \DisplayProof
      &
      \AxiomC{\sFmla{\False}{\Diamond !A}[n]}
      \RightLabel{\TRule{\False}{\Diamond}}
      \UnaryInfC{\sFmla{\False}{!A}[m]}
      \DisplayProof\\
      $m$ جديدة & $m$ مستعملة
      \\[1ex]
      \hline}{}
    \end{tabular}
  \end{center}
  \caption{قواعد مبسطة للنسق \Log{S5}.}
  \ollabel{tab:rules-S5}
\end{table}

\begin{ex}
  نعطي جدولًا دلاليًا مغلقًا مبسطًا يبيّن أن $\Log{S5} \Proves
  \Ax{5}$، أي $\Diamond!A \lif \Box\Diamond!A$.
  \begin{oltableau}
    [\pFmla{\False}{\Diamond\formula{A} \lif \Box\Diamond \formula{A}}{1},
      just = \TAss
      [\pFmla{\True}{\Diamond \formula{A}}{1}, just = {\TRule{\False}{\lif}[1]}
        [\pFmla{\False}{\Box\Diamond \formula{A}}{1},
          just = {\TRule{\False}{\lif}[1]}
          [\pFmla{\False}{\Diamond \formula{A}}{2},
            just = {\TRule{\False}{\Box}[3]}
            [\pFmla{\True}{\formula{A}}{3},
              just = {\TRule{\True}{\Diamond}[2]}
                [\pFmla{\False}{\formula{A}}{3},
                  just = {\TRule{\False}{\Diamond}[4]}, close]
            ]
          ]
        ]
      ]
    ]
  \end{oltableau}
\end{ex}

\end{document}
